\centerline{\bf \Huge Тренировка региона ВсОШ}
\centerline{\bf \Huge 10 класс}

\problem{10.1}
Первоклассник составил из шести палочек два треугольника. Затем он разобрал треугольники обратно и разбил шесть палочек на две группы по три палочки: в первой группе оказались три самых длинных палочки, а во второй - три самых коротких. Обязательно ли можно составить треугольник из трех палочек первой группы? А из трех палочек второй группы?

\problem{10.2}
Ненулевые числа $x$ и $y$ удовлетворяют неравенствам $x^4-y^4>x$ и $y^4-x^4>y$. Может ли произведение $x y$ равняться отрицательному числу?

\problem{10.3}
Пусть $S$ - множество, состоящее из натуральных чисел. Оказалось, что для любого числа $a$ из множества $S$ существуют два числа $b$ и $c$ из множества $S$ такие, что $a=\frac{b(3 c-5)}{15}$. Докажите, что множество $S$ бесконечно.

\problem{10.4}
Вписанная окружность касается сторон $A B, B C$ и $C A$ неравнобедренного треугольника $A B C$ в точках $C_1, A_1$ и $B_1$ соответственно. Пусть $m$ - средняя линия треугольника $A_1 B_1 C_1$, параллельная стороне $B_1 C_1$. Биссектриса угла $B_1 A_1 C_1$ пересекает $m$ в точке $K$. Докажите, что описанная окружность треугольника $B C K$ касается $m$.

\problem{10.5}
Петя и Вася играют на доске $100 \times 100$. Изначально все клетки доски белые. Каждым своим ходом Петя красит в чёрный цвет одну или несколько белых клеток, стоящих подряд по диагонали. Каждым своим ходом Вася красит в черный цвет одну или несколько белых клеток, стоящих подряд по вертикали. (На рисунке справа показаны возможные первые ходы Пети и Васи на доске $4 \times 4$.) Первый ход делает Петя. Проигрывает тот, кто не может сделать ход. Кто выигрывает при правильной игре?